\chapter{Convergence of MCMC and the Witch's Hat}
\label{ch:witchhat}

Every level set sampler in this book is a Markov chain Monte Carlo algorithm, and the question of its mixing is the question statisticians have asked about Gibbs and Metropolis samplers since the late 1980s. This chapter gives the statistical side of that story: the Gibbs sampler and the Metropolis Hastings algorithm as they are used in Bayesian inference, the Peskun ordering, the convergence theory of Tierney and of Meyn and Tweedie, the diagnostics of Gelman and Rubin and of Geyer, and the witch's hat example which showed that a chain can look converged while being exponentially far from it. The witch's hat is, in miniature, the problem of the bottom rung.

\section{Gibbs and Metropolis Hastings}

The Gibbs sampler, named by Geman and Geman in 1984 and brought into mainstream statistics by Gelfand and Smith in 1990, updates a vector $x = (x_1, \dots, x_d)$ one coordinate at a time, drawing $x_i$ from its conditional distribution given the rest. Each update leaves the target invariant, and the chain is reversible if the coordinates are updated in random order. The level set chain for the queens is a Metropolis chain on pairs of coordinates: a swap of rows $i$ and $j$ proposes the other order of $(x_i, x_j)$ and accepts if the energy stays below the threshold. The Gibbs sampler on the pair, which draws $(x_i, x_j)$ from its conditional given the other rows, uniform on the at most two admissible values, is the lazy version of this chain: it swaps with probability $1/2$ when both orders are admissible.

The Metropolis Hastings algorithm, from Metropolis and coauthors in 1953 and Hastings in 1970, proposes $y \sim q(x, \cdot)$ and accepts with probability
\[
\alpha(x, y) = \min\Big\{ 1, \, \frac{\pi(y) \, q(y, x)}{\pi(x) \, q(x, y)} \Big\} ,
\]
which makes the chain reversible with respect to $\pi$ for any proposal $q$. Gibbs is the special case in which the proposal is the conditional and the acceptance is always one. The block chain of Chapter~\ref{ch:swendsen} is Metropolis Hastings with a block proposal, and the generator proposes and chain corrects loop of Chapter~\ref{ch:gbc} is Metropolis Hastings with a learned proposal.

\section{Peskun's ordering}

Peskun showed in 1973 that among reversible chains with the same stationary distribution, one that moves more, in the sense that its off diagonal transition probabilities dominate those of another, has smaller asymptotic variance for every function. The Metropolis acceptance rule is optimal in Peskun's order among rules of the form $\alpha(x, y) = s(x, y) / (1 + t(x, y))$ with $s$ symmetric and $t(x, y) = \pi(x) q(x, y) / (\pi(y) q(y, x))$, which is why it is used rather than the Barker rule $\pi(y) / (\pi(x) + \pi(y))$. For the queens the swap chain on a level set accepts every proposal that stays in the level set and rejects every proposal that leaves it, which is the Metropolis rule for a uniform target, and no rule in this family does better. The block chain dominates the swap chain in Peskun's order when its transition probabilities satisfy $P_{\rm block}(x, y) \ge P_{\rm swap}(x, y)$ for all $x \ne y$, and then has no larger asymptotic variance, though its cost per step is higher; containing the swap proposal as a special case does not by itself give this ordering.

\section{Convergence theory}

Tierney's 1994 paper gave the conditions under which a Metropolis Hastings chain converges: irreducibility, aperiodicity and Harris recurrence suffice for the law of large numbers, and geometric ergodicity, which is a drift and minorisation condition in the sense of Meyn and Tweedie, gives the central limit theorem with a finite asymptotic variance. One of us, in a 1996 discussion paper, developed these conditions for Gibbs and Metropolis chains as used in Bayesian computation and gave explicit drift functions for several standard models, and argued that the rate of convergence of a Bayesian sampler is governed by the same spectral quantities, gap and conductance, as the chains of theoretical computer science. The argument was that the two literatures had been studying one object under two names, and the present book is that argument carried through.

Geometric ergodicity says that the distance to stationarity decays like $\lambda^t$ for some $\lambda < 1$. It does not say what $\lambda$ is, and for a chain on an exponentially large state space $\lambda$ may be exponentially close to one. The queens swap chain on a level set is geometrically ergodic whenever it is irreducible and aperiodic (for instance lazy), trivially, since the state space is finite; the question of its mixing time is the question of whether $1 - \lambda$ is inverse polynomial, and geometric ergodicity alone does not answer it.

\section{Diagnostics}

In practice one monitors a chain by diagnostics. Gelman and Rubin's $\hat R$ compares the between chain and within chain variances of several chains started from dispersed points, and signals convergence when the ratio is near one. Geyer's initial sequence estimators bound the asymptotic variance from the autocorrelations of a single chain. The effective sample size, $n / \tau_{\rm int}$ for a chain of length $n$ and integrated autocorrelation time $\tau_{\rm int}$ (not the mixing time $\tau$ of the Notation chapter), is built from the integrated autocorrelation time we reported in Chapter~\ref{ch:experiments2}. All of these can detect a failure to converge and none can certify convergence: a chain trapped in a region of the state space will have autocorrelations that decay within the region and chains started within the region will agree, and the diagnostics will report convergence to a distribution that is not the target.

\section{The witch's hat}

Geyer's 1992 paper on practical Markov chain Monte Carlo argued for long single runs over many short ones, and one of us gave in the discussion an example, the witch's hat, that shows what the diagnostics cannot see. The target density on the unit cube in $d$ dimensions is a mixture of a broad uniform component and a narrow spike, a peaked cone of small width centred at an interior point, with the spike carrying a fixed fraction of the mass. The Gibbs sampler on this target updates one coordinate at a time from its conditional, which is nearly uniform unless the other coordinates already lie within the spike's width of the centre, and the probability of that happening by chance is the spike's width to the power $d - 1$, exponentially small. A chain started in the brim of the hat stays in the brim for a time exponential in $d$; a chain started in the spike stays in the spike. Each chain, examined by any diagnostic, looks perfectly converged, with autocorrelations decaying in a few steps and $\hat R$ near one, and each is sampling from the wrong distribution, missing a constant fraction of the mass.

The witch's hat is the simplest example of a target with a tiny basin reached only through an exponentially small window: the set of points of density above the brim's level is the spike alone, and the chain's local moves almost never land in it. The example made the point in 1992 that convergence diagnostics cannot detect multimodality that the chain has not yet encountered, and the point is exactly the one that Chapter~\ref{ch:mixing} makes about the queens: a swap chain on a level set with several components looks converged within its component. The exponential time to find the spike plays the role that the exponential mixing time of a chain across components would play, and for the queens on hard boards we expect it to be the obstruction to $\NP = \RP$.

\section{From the witch's hat to the queens}

The witch's hat has a single spike at a known location; a hard queens board has, we expect, exponentially many spikes at unknown locations, the components of the level set, and a chain that finds one of them has no way to know how many it has missed. The genealogical diagnostic of Chapter~\ref{ch:smc}, the number of distinct ancestors of the particles that reach the bottom, is a diagnostic that can see part of this: a particle system started from many dispersed points in the brim will, if the spikes are few and findable, have ancestors in several of them, and if they are many and hidden will coalesce to the one it happened to find, so the diagnostic reveals the missed spikes only when the dispersed starting points reach them. It is the multichain diagnostic of Gelman and Rubin applied to the bottom rung with the chains coupled through selection, and it is the right tool because it measures diversity where diversity is lost.

The lesson of the witch's hat for the program is that the Monte Carlo evidence of Chapters~\ref{ch:experiments} and~\ref{ch:experiments2}, which showed the chain working on unconditioned boards and boards conditioned on queens of a solution, is evidence about the components the chain found and not about the components it did not. The exact connectivity computations of Chapter~\ref{ch:open} are what rule out hidden components at small $N$, by enumeration rather than sampling, and they are the only kind of evidence that can. At large $N$ on adversarial boards there is no enumeration, and the question is open in the strongest sense: not merely unproved but, as the witch's hat illustrates, possibly undetectable by chain based diagnostics.
